\documentclass[11pt, letterpaper]{article}
\input{../../homework.tex}
\usepackage{titlepageBU}
\title{Turn in 4}
\courseID{MA 542}
\courseSection{A1}
\begin{document}
\makeproblem
Note: there is a lot of explicit matrix computation involved in this assignment. I have included all computations at the end just incase you want to see them, but I wasn't sure and didn't want to make it too messy so like yeah.

For simplicity, denote
\[
	M_1\coloneqq \begin{pmatrix}
	1 & 0 \\
	0 & 0
	\end{pmatrix},\quad M_2\coloneqq \begin{pmatrix}
	0 & 1 \\
	0 & 0
	\end{pmatrix},\quad M_3\coloneqq \begin{pmatrix}
	0 & 0 \\
	1 & 0
	\end{pmatrix},\quad M_4\coloneqq \begin{pmatrix}
	0 & 0 \\
	0 & 1
	\end{pmatrix}
.\]
Also let $1_R$ denote the identity matrix. Explicit computation yields that
\[
	M_1M_2=M_2,\quad M_3M_2=M_4,\quad M_4M_3=M_3,\quad M_2M_3=M_1
.\]
This can be visualized with the diagram:
\[\begin{tikzcd}[row sep = 2.5em, column sep = 2.5em]
	&M_2\arrow[dr, bend left, "M_3"]\\
	M_1\arrow[ur, bend left, "M_2"]&&M_4\arrow[dl, bend left, "M_3"]\\
				       &M_3\arrow[ul, bend left, "M_2"]
\end{tikzcd}\]
In some cases the operation is left-multiplication, and in others the operation is right-multiplication. The important thing to note is that if there exists some $i\in \left\{ 1, \ldots , 4 \right\} $ with $M_i\in I$, then it follows that $M_i\in I$ for all $i$, because $I$ is an ideal and thus has the absorbtion property, and all $M_i$ can be represented as a product in $R$ on a single $M_i$ by the above diagram. I will now show the ideal contains at least one of these.

Let
\[
	A=\begin{pmatrix}
	a & b \\
	c & d
	\end{pmatrix}\in I
\]
be nonzero. Note that $A=aM_1+bM_2+cM_3+dM_4$. Since $A$ is nonzero, at least one element in $\left\{ a,b,c,d \right\} $ must be nonzero. If precisely one element is nonzero, then we have some coefficient $k\in\mathbb{R} $ such that
\[
	A=kM_i
\]
for some $i$. Since $I$ is an ideal, $A\in I$, and $\frac{1}{k} 1_R\in R$, we have $\frac{1}{k} 1_R A=\frac{k}{k} 1_RM_i=M_i$ is an element of the ideal.

Now suppose that more than one coefficient is nonzero. In particular, we must have at least two coefficients are nonzero. Consider the following three cases:
\begin{itemize}
	\item \textbf{Case 1:} $a,b\neq 0$, and $c,d=0$. That is,
		\[
			A=aM_1+bM_2
		.\]
		Explicit computation shows that $M_1M_4=0$ and $M_2M_4=M_2$. It follows that
		\[
			AM_4=\left( aM_1+bM_2 \right)M_4 =aM_1M_4+bM_2M_4=bM_2
		.\]
		Here, I use some facts about vector spaces, such as $a0_R=0_{R}$ for all scalars $a$. Since $I$ is an ideal and $a\in I$, we have $AM_4=bM_2\in I$. By the same logic, we also have
		\[
			\frac{1}{b} 1_R\left( bM_2 \right) =1_RM_2=M_2\in I
		.\]
	\item \textbf{Case 2:} $c,d\neq 0$, and $a,b=0$. That is,
		\[
			A=cM_3+dM_4
		.\]
		Explicit computation yields that $M_3M_1=M_3$ and $M_4M_1=0$. Therefore, we have
		\[
			AM_1=\left( cM_3+dM_4 \right) M_1=cM_3M_1+dM_4M_1=cM_3
		.\]
		Since $A\in I$ and $I$ is an ideal, $AM_1=cM_3\in I$. Similarly,
		\[
			\frac{1}{c} 1_R\left( cM_3 \right) =M_3\in I
		.\]
	\item \textbf{Case 3:} Any other possible scenario. We must have either $a$, $b$, or both are nonzero. By explicit computation, we have
		\[
			M_1M_1=M_1,\quad M_1M_2=M_2,\quad M_1M_3,M_1M_4=0
		.\]
Since left multiplication fixes $M_1,M_2$ and annihilates $M_3,M_4$, we can take $c,d=0$. If only one of the coefficients $a,b$ are nonzero, then we have the corresponding $M_i\in I$ by previous logic. If both are nonzero, apply the same logic as in case 1.
\end{itemize}
These are the only cases, since $A\neq 0$. Therefore, $M_1,M_2,M_3,M_4\in I$.

Since $I$ is an ideal, it must be a subring, and thus closed under addition. We thus have
\[
	M_1+M_4=\begin{pmatrix}
	1 & 0 \\
	0 & 0
	\end{pmatrix}+\begin{pmatrix}
	0 & 0 \\
	0 & 1
	\end{pmatrix}=1_R\in I
.\]
Let $M\in \mathcal{M}_2(\mathbb{R} )$ be any matrix. Then since $1_R\in I$, and $I$ is an ideal, we have that $1_RM=M\in I$. Therefore, $\mathcal{M}_2(\mathbb{R} )\subseteq I$. Since $I$ is an ideal, we have $I\subseteq \mathcal{M}_2(\mathbb{R} )$, and thus $I=\mathcal{M}_2(\mathbb{R} )$.

\[
	i \hbar \frac{\mathrm{d}}{\mathrm{d}t} \ket{\Psi (t)} =\mathbf{\hat{H} } \ket{\Psi (t)} 
.\]

\end{document}
this can totally be reduced to two cases. if a != 0 or b != 0, then M1 A collapses. if only one, we're done, so assume 2. then take M1 A M3 = aM2. Also if a,b = 0 then M2 A = cM1 + dM2 so were back to the same thing as before

M1M1=M1
M1M2=M2
M1M3,4=0

M2M1,2=0
M2M3=M1
M2M4=M2

M3M1=M3
M3M2=M4
M3M3,4=0

M4M1,2=0
M4M3=M3
M4M4=M4
